% =====================================================================================================
\section{Discussion and Limitations}\label{sec:discussion}
\textit{Limitations.} First, power is not measured: the energy per frame combines measured frame rates with Vivado's estimate. On the board a complete frame is set by the DMA transfers and the ARM's headers, not by the encoder, and our driver starts the DMA after the encoder, which adds a fixed wait to the fabric time. Second, sign and amplitude encryption leak outlines by construction (Lemma~\ref{lem:format}); only a change of block positions removes them, at a byte cost. Third, frame sizes leak at every level (Section~\ref{sec:sizes}). Fourth, our attackers are greedy or trained on a few clips; stronger solvers or attackers trained on large corpora may extract more from E5, and a trained network already undoes sign-only encryption. Fifth, a pre-shared key must still be provisioned, but it only authenticates the hybrid ML-KEM-768~\cite{fips203} + X25519 exchange that yields each session's key, much as DTLS-SRTP keys SRTP~\cite{rfc3711}; a recorded session stays secret if the pre-shared key leaks later or a quantum computer breaks the elliptic curve, and the exchange takes \BdKexMs\,ms on the Cortex-A9 with a pure-Python ML-KEM, once per session. The fabric needs no change: Ascon-AEAD128 is symmetric, and Grover's search~\cite{grover1996} needs on the order of $2^{64}$ evaluations for a 128-bit key, a reduction of the margin, not a break. Physical attacks would call for a threshold implementation of Ascon~\cite{gross2017}.

\textit{Reproducibility and ethics.} Every number in this article is produced by a script from stored results and inserted by a generator, never typed; a test fails if a table differs from the data. The golden model, the RTL, the testbenches, the Vivado and board scripts, the board images and the experiment scripts are available at \url{https://github.com/bnsama29-cloud/FPGA-Video-Encryption-}. The private clips are identified by hash; no picture derived from them appears here or is stored.

% =====================================================================================================
\section{Conclusion}\label{sec:conclusion}
Encrypting coefficient signs and amplitudes with Ascon inside a JPEG encoder costs no cycles, provably keeps every protection level playable and, with a per-frame tag over metadata and scan, turns a plain MJPEG stream into an authenticated one. Field encryption alone leaves the outline of the scene to anyone who counts non-zero coefficients, and published format-compliant schemes fall to one known frame because their key material does not change. A keyed permutation of macroblocks drawn anew for every frame removes the outline under every attack we could build, and with restart segments it fits in hardware without a frame buffer. On a PYNQ-Z2 board every image is bit-exact, four lanes encode \BdFourBusyMpix\,Mpixel/s at \BdFourMhz\,MHz, a complete authenticated frame takes \BdFourZcMs\,ms, and a closed feedback loop mitigated every attack in simulation and with the board as the camera. Next are measured power, a driver that starts the DMA before the encoder, size padding in hardware and a shared Ascon core across lanes.
